% 第 58 章　诺特定理：守恒律从哪里来
\chapter{诺特定理：守恒律从哪里来}

第~57~章把“对称”从一句夸奖变成了一个可以操作的概念：对系统做某种改变，如果物理规律纹丝不动，就说规律具有这种对称。本章要解决一个更大的问题。全书的账本体系——第~9~章的能量、第~10~章的动量、第~11~章的角动量——一直建立在三条“守恒律”之上：孤立系统的总能量、总动量、总角动量各自保持不变。我们一直像会计一样用它们对账，对上了就安心，对不上就去找藏起来的那笔。可是你有没有追问过：这三条“禁令”是谁颁布的？为什么恰好是这三个量守恒，而不是速度、加速度或者别的什么东西？一九一八年，数学家埃米·诺特给出了答案，而且是一个可以证明的答案：守恒律不是额外的禁令，它们是对称性的影子。本章就来看这个影子是怎么投出来的。

\step{反常识开场}

花样滑冰的比赛录像里，有一个镜头值得倒回去反复看。运动员单脚立在冰面上开始旋转：起初双臂平伸，像一架缓慢的风车，大约一秒转一圈；接着她把双臂收拢抱向胸口——没有蹬地，没有人推，转速却猛地蹿上去，快到面目模糊，一秒能转三四圈。然后她重新张开手臂，转速又乖乖地慢下来，为下一个动作腾出从容的时间。

慢下来想一想这件事有多奇怪。第一，她变快不需要任何“推力”。冰面几乎不提供水平方向的力，竖直方向的支持力和重力都与转轴平行——谁也没有拧她一把，可速度的增加就像凭空变出来的。第二，更蹊跷的是这件事精确得可怕。测量运动员的质量分布，可以算出她的转动惯量——质量离转轴越远，这个量越大。把转速和转动惯量的乘积记下来：收臂前是某个数，收臂后还是同一个数。转动惯量缩到三分之一，转速就恰好涨到三倍，不多不少，仿佛冰面上有一台看不见的收银机，每笔账都分毫不差。

再看另一件你更可能亲身经历过的事。从一条静止的小船往岸上跳：你向前跃出的一瞬，船必然向后滑去。停下来量一量会发现，你的质量乘上你的速度，和船的质量乘上船的速度，大小几乎正好相等，方向正好相反。跳得猛，船退得快；跳得轻，船退得慢；两个人配合着跳也一样——账永远平。

第三件事不起眼到你可能从没注意过它：今天重复昨天的实验，结果是一样的。你家楼下的秋千，昨天荡一个来回要几秒，今天还是要几秒；厨房里的小苏打遇到醋，昨天冒泡，今天还冒泡。你从来没有担心过“物理规律隔夜涨价”，写实验报告时默认仪器的常数可以沿用一整个学期。这种心安理得本身，就是一种惊人的宇宙行为。顺带一提，你此刻正随地球以每秒约三十公里的速度绕太阳狂奔，又随太阳系以每秒二百多公里的速度绕银河系中心漂移——可是你跳绳的成绩、杯中茶水的温度，与这些浩浩荡荡的速度毫无关系。如果物理规律在乎“你在宇宙里走多快”，你的生活早就乱套了。

三件事摆在一起，就是本章的悬念。物理学里有三条著名的“禁令”：孤立系统的总动量不变、总角动量不变、总能量不变。运动员的收银机对应角动量守恒，跳船的对称账目对应动量守恒，而“实验隔夜不变”的背后站着能量守恒。禁令人人会背，可禁令从哪里来？为什么它恰好管着这三个量，而不管速度、不管加速度？

\step{先猜一猜}

把你的直觉立场先写下来，越具体越好。

第一题。关于守恒律的来历，有两种立场。甲说：守恒律是宇宙额外颁布的三条禁令，和运动定律并列，互不隶属，不需要更深的理由；乙说：守恒律不是独立的法律，而是某种更深、更平常的事实的影子。你的直觉站哪一边？

第二题。设想三次“搬家”：把一整套实验设备从你家搬到一千公里外的另一座城市重做一遍；把整套设备在原地转过九十度重做一遍；把整套实验原封不动地留到明天重做一遍。你猜，三次搬家的结果会不会变？如果某一次搬家之后结果真的变了，你觉得守恒律会出什么事？

第三题。推一个箱子在地面上滑，箱子很快停下来——动能没了。有人会说：“能量守恒只是会计技巧，账不平了就开一个新科目叫‘热’，账就永远平。”你猜：这个“新科目”是真的对应着某种可以独立测量的东西，还是仅仅是账房先生的遮羞布？

写好后夹在书页里，本章结束时逐条对答案。

\step{动手实验}

\begin{experiment}[转椅上的收银机]
\textbf{材料}：一把能自由转动的电脑转椅，两本厚书或两瓶装满水的矿泉水。

\textbf{第一步}：坐进椅子，双脚离地，双臂向两侧平举，每只手握一本书。请同伴轻轻推你的手臂，让椅子缓慢转起来，大约两三秒一圈即可，不要快。

\textbf{第二步}：转动过程中，把双臂收拢抱到胸前。你会清楚地感到转速一下子变快——没有任何人推，只是收臂。

\textbf{第三步}：再把双臂张开，转速又慢回去。如此反复几次，体会“快慢由手臂说了算”的感觉。

\textbf{第四步（选做）}：转动中把上半身明显地向一侧倾斜，你会感到椅子的转动变得“别扭”、晃荡，不再绕着原来的竖直轴干净地转。

\textbf{安全提醒}：始终慢速，收臂时握紧书本，第一次做请同伴在旁扶稳。

\textbf{记录与思考}：让同伴帮你数“十秒几圈”，估一估收臂前后转速大约变了几倍。然后回答两个问题：转速明明变了，你却感觉不到任何人在“拧”椅子——变的那个量和不变的那笔“账”，各是什么？第四步里，转动为什么变得别扭？
\end{experiment}

这个实验还藏着一个更深的谜团，先记下来：收臂之后你转得更快，转动动能反而增加了——能量好像也不守恒了？账房先生是不是又开了新科目？本章后面会把这笔账也算平。

\step{旧模型遇到麻烦}

先审问第一种旧模型：“守恒律就是宇宙颁布的三条禁令，与别的定律并列。”它的麻烦在于解释不了“为什么偏偏是这三个量”。为什么守恒的是动量而不是速度？是能量而不是某种“力的累计”？禁令的清单为什么不多不少（在三维空间里）正好三条？如果禁令是任意颁布的，那么清单长成什么样都没有道理可讲——可它偏偏严丝合缝。

再审问第二种旧模型：“守恒律是从牛顿运动定律推出来的副产品。”逻辑上确实可以推：用牛顿第三定律，两个物体相互作用时动量此消彼长，总动量不变。这个推导本身没错，但后来的历史证明它把因果弄反了。电磁理论发现光也携带动量——太阳帆能被光推着慢慢加速，光压实验里镜子获得的动量恰好等于光动量的变化。牛顿力学里根本没有“场的动量”这个位置，可动量守恒照样成立。再往后，量子力学干脆连“力”这个词都很少使用，可能量、动量、角动量守恒在量子世界不但成立，还反过来当上了立法者：一个核反应、一次粒子衰变能不能发生，先查账平不平。如果守恒律只是牛顿定律的副产品，它应当随着牛顿定律的失效而失效；事实是，牛顿定律在微观和高速领域退了场，守恒律却越活越精神，管辖范围反而更大了。

还有一个更日常的尴尬，你刚在转椅上已经撞见过：收臂之后动能增加了。刹车时汽车的动能变成刹车片的滚烫，两个泥球对撞粘在一起后动能几乎全灭。如果能量守恒是铁律，这些“消失”和“冒出”的能量去了哪、从哪来？旧模型要么含糊其辞，要么临时开科目。而“禁令模型”的回答永远是“因为这是禁令”——这不是解释，是复读。

二十世纪初的物理学已经攒了一堆这样的疑问。一九一五年，爱因斯坦完成广义相对论之后，数学家们被请去帮忙清理其中的能量守恒问题；一九一八年，德国数学家埃米·诺特给出了回答。她的回答不是一条新的禁令，而是一个定理：守恒律的清单不是任意的——物理规律每具有一种连续的对称，就必然存在一个对应的守恒量。禁令是影子，对称才是物体。

\step{发明新概念}

先回忆第~57~章的定义：对称就是“做了某种改变之后，规律不变”。本章需要三种连续对称——“连续”的意思是改变可以一点点地做，而不像镜面反射那样只有“换”与“不换”两档。它们是：空间平移，把整套实验挪个地方；时间平移，把整套实验推迟一段时间；空间旋转，把整套实验转个方向。现在逐个看看，这三种“不变”各自会变出什么守恒量。

\textbf{空间平移不变，对应动量守恒。}回到冰面上互推的两个人。如果规律处处相同，那么“甲推乙”所遵循的规则和“乙推甲”所遵循的规则必须是同一套规则，只是换了主语：两个人的位置没有任何特殊之处，任何一处都不比另一处更“配”被规律优待。反过来，假如规律真的随地点变化，就没有任何东西保证这两份作用成对抵消，孤立系统的总动量就可以凭空增减。换句话说：总动量会改变，当且仅当规律本身随地点改变。规律处处相同，动量的账就被迫是平的。

\textbf{时间平移不变，对应能量守恒。}这里用一个“奸商论证”。假设物理规律随时间变化，比如引力在夜里变弱、白天变强。那么你可以夜里用很少的功把水抽上山顶的水库，白天再放水发电，发出的电比抽水耗的多；循环操作，净赚能量，永动机就此合法化。能量守恒禁止一切永动机，等价于说：不存在“夜里更便宜”这种时间上的特价——规律不随时间变。能量守恒，就是“规律没有时差”这句话的记账形式。

\textbf{空间旋转不变，对应角动量守恒。}回到转椅。如果规律对各个方向一视同仁，就不存在任何“特权方向”可以让转动趋势凭空漏走或凭空注入；绕某根轴的转动账只能自己摆平。反过来，如果真的出现了特权方向——你斜身子那一下，引力替你选了“下”这个方向——绕别的轴的角动量就可以改变，转动立刻变得别扭。

三段论证有一副共同的骨架：某个量守恒，当且仅当规律对某种改变“无所谓”。诺特定理把这副骨架变成了严格的数学：只要运动规律可以由一个“作用量”（第~13~章那位老朋友）写出来，那么作用量具有的每一种连续对称，都精确对应一个沿真实运动保持不变的数量。这是定理，不是猜想——给出对称，守恒量可以被公式直接算出来。更有意思的是，定理管的不只是“有没有”，还管“长什么样”：守恒量的具体形式（动量是质量乘速度、角动量是转动惯量乘角速度）同样由对称的数学结构算出，而不是额外的假设。

请注意“连续”两个字，它是定理的门槛。镜面反射、左右互换这类离散对称，只有“换”与“不换”两档，没有“换一点点”的中间态，诺特的推导机器对它们使不上劲——离散对称不对应守恒量。这不等于说离散对称没用：镜像对称在量子力学里给出“宇称”这个好用的标签，能判断哪些过程允许、哪些禁戒；但它给出的不是一个可以连续取值的守恒数量。一九五六年之前，物理学家曾默认镜像对称和左右旋守恒一样天经地义，直到实验发现弱相互作用居然分得清左右——这是第~59~章和第~60~章那片区域的风景，这里只记住一句话：对称的家谱里，只有连续的那一支，才派发守恒量。

现在可以顺手解开旧模型留下的那个疙瘩：禁令清单为什么“不多不少正好这些”？数一数对称的个数就知道了。三维空间里，独立的平移方向有三个（前后、左右、上下），对应动量的三个分量各自守恒；独立的转动轴有三根，对应角动量的三个分量；时间只有一个方向，所以能量只有一个。三维空间里的“禁令”不是三条，而是三加三加一共七条——每一条都挂在一根具体的对称轴上。假如有一天人类发现规律对某种全新的连续改动也无所谓，守恒量的清单上就会多出一条；反过来，清单上没有任何一条是多余的，因为背后没有闲着的对称。清单的形状，就是时空对称性的形状。

作用量的直观值得再花一段。把作用量想成整条路线的“总票价”：像票价的地方在于，要比较的是全程的总账，而不是某一步的贵贱；不像票价的地方在于，作用量有它自己的物理单位（能量乘时间），而且自然只要求“驻值”——把路线改动一点点，总账的一阶变化为零——并不一定真的是最小。现在把对称放进来：对称的意思就是“总票价对某种改动完全不敏感”，把整条路径平移一下、推迟一下、旋转一下，总票价原封不动。而“沿真实路径走时，总账对某个方向不敏感”在数学上恰好等价于“某个数量沿路径保持不变”。不敏感的方向，就是守恒的方向。这就是诺特定理的全部直觉。

\begin{figure}[htbp]
\centering
\begin{tikzpicture}[>=Stealth,
  box/.style={draw,rounded corners=3pt,align=center,minimum width=4.6cm,minimum height=1.0cm,fill=blue!4},
  cbox/.style={draw,rounded corners=3pt,align=center,minimum width=4.6cm,minimum height=1.0cm,fill=green!6}]
  \node[box] (s1) at (0,0) {空间平移不变\\（在哪里做都一样）};
  \node[box] (s2) at (0,-1.7) {时间平移不变\\（什么时候做都一样）};
  \node[box] (s3) at (0,-3.4) {空间旋转不变\\（朝哪个方向做都一样）};
  \node[cbox] (c1) at (8,0) {动量守恒};
  \node[cbox] (c2) at (8,-1.7) {能量守恒};
  \node[cbox] (c3) at (8,-3.4) {角动量守恒};
  \draw[->,thick] (s1) -- (c1);
  \draw[->,thick] (s2) -- (c2);
  \draw[->,thick] (s3) -- (c3);
  \node at (4,0.75) {规律“无所谓”的方向};
  \node at (4,-4.15) {被迫保持不变的数量};
  \draw[->,thick,dashed] (2.6,0.45) .. controls (4,0.3) .. (5.4,0.45);
\end{tikzpicture}
\caption{诺特定理的三对“对称—守恒”：左边是规律不敏感的三种改变，右边是被迫保持不变的三个数量。箭头是定理保证的严格对应，不是类比。}
\end{figure}

补一条历史注脚。埃米·诺特是二十世纪最伟大的数学家之一，却因性别长期拿不到正式教职，在哥廷根只能以别人的名义开课。她的这条定理发表之初反响平平，几十年后却成了理论物理的底层语法：今天，粒子物理学家写下任何一个新理论，第一件事就是清点它的对称性，因为对称性决定了守恒量，而守恒量决定了这个理论允许什么、禁止什么。一个连职称都拿不到的人，给整个宇宙立了规矩。

\step{一句话模型}

\begin{oneline}
宇宙并没有颁布“动量不许变、能量不许变、角动量不许变”这三条禁令；它只是在三个方向上“无所谓”——规律不因为你换了地方、换了时刻、换了朝向而改变。规律对平移无所谓，孤立系统的总动量就被迫不变；对时间推移无所谓，总能量就被迫不变；对转向无所谓，总角动量就被迫不变。守恒律是对称性的影子：影子不是额外的物体，而是物体挡光的方式。
\end{oneline}

\step{数学翻译器}

先看动量。平移对称逼出“作用与反作用等大反向”：\(\mathbf F_{12}=-\mathbf F_{21}\)。合力改变动量（第~6~章和第~10~章的约定），于是
\[
\frac{d}{dt}\bigl(\mathbf p_1+\mathbf p_2\bigr)=\mathbf F_{21}+\mathbf F_{12}=0 .
\]
按老规矩回答四个问题。描述什么：两个物体组成的孤立系统，其总动量随时间的变化率——等于零意味着总动量是一个常数。为什么这样写：左边是“总账的变动”，右边两项是系统内部仅有的两笔内力，平移对称保证它们互相抵消；式子把“规律处处相同”翻译成了“总账变动为零”。何时成立：系统孤立，或者外力之和为零——冰面上互推的两人，竖直方向重力与支持力抵消，水平方向近似无摩擦。何时失效：存在净外力时——在地面上互推，摩擦力把动量传给了地球，两人的总动量不再守恒；但“两人加地球”的总动量依然守恒。失效不是守恒律破了，而是你把账本的边界画小了。

用特殊情形检查。总动量为零：静止的两人互推之后，动量一正一负之和仍为零——跳船的人和船正是此例：\(60\)~kg 的人以 \(2\)~m/s 上岸，\(40\)~kg 的船以 \(3\)~m/s 后退，两笔乘积都是 \(120\)~kg\(\cdot\)m/s，方向相反。方向反转：两颗相同的台球正面对撞后交换速度，式子不管碰撞过程多复杂，只管首尾两笔账。极大情形：超新星爆发把星体炸向四面八方，所有碎片的总动量仍等于爆发前那颗恒星近乎为零的动量——天文学家当年就是靠这笔账预言了中微子：可见碎片加上辐射的账对不平，必定有“隐身碎片”把动量带走了，后来果然被逮到。

再看角动量。绕某根轴的转动惯量 \(I\) 与角速度 \(\omega\) 的乘积记为 \(L=I\omega\)，旋转对称给出
\[
L=\text{常量}\quad\text{（绕对称轴、无外力矩）}.
\]
描述什么：转动趋势的“总额”。为什么这样写：\(I\) 管质量离轴有多远，\(\omega\) 管转多快；收臂改变前者，后者必须跟着变，才能保住乘积。何时成立：规律绕该轴旋转对称，且没有外力矩。何时失效：引力选了“下”这个方向时，只有绕竖直轴的分量守恒——你在转椅上歪身子，账就乱了。代入花滑的真实数字：伸臂时 \(I\approx3.0\)~kg\(\cdot\)m\(^2\)，收臂后 \(I\approx1.0\)~kg\(\cdot\)m\(^2\)，转速从约每秒一圈涨到约每秒三圈，乘积分毫不差。再看动能 \(E=\tfrac12 I\omega^2=L^2/(2I)\)：角动量不变而转动惯量缩到三分之一，动能反而涨到三倍——从约 \(59\)~J 涨到约 \(178\)~J。多出来的能量不是凭空来的：收臂时肌肉要把手臂从甩出去的趋势中“拽”回来，做了正功，账从化学能科目转入动能科目。角动量守恒从不承诺动能守恒。再查极端：\(I\) 不变则 \(\omega\) 不变——空竹不推不挡就匀速空转（忽略摩擦）；\(\omega=0\) 则 \(L=0\)，怎么收臂也转不起来，零乘任何数还是零——这解释了你为什么必须先被推一下。

这条看似只能解释冰上舞蹈的定律，其实是现代导航的根基。一只高速旋转的陀螺，角动量的方向被旋转对称“锁”在空间里：载体怎么颠簸转弯，转轴方向都稳稳不动。把三只陀螺按互相垂直的方向架起来，就得到一套不依赖任何外界信号的“方向记忆”——飞机、潜艇和洲际导弹的惯性导航靠它，你的手机能感知屏幕横竖和步数，靠的也是指甲盖大小的微机械陀螺。反作用轮与陀螺仪，一个是“主动调账”，一个是“被动守账”，用的是同一条对称。

时间平移翻译成公式稍微绕一点，主线只记结论：只要系统的规律不随时间变化——用第~15~章的语言说，哈密顿量不显含时间——就有一个叫“总能量”的数量保持不变：
\[
E=\text{常量}\quad\text{（规律不随时间变）}.
\]
这里要诚实交代一个前提，第~15~章埋过伏笔：“哈密顿量守恒”和“哈密顿量等于动能加势能那份总能量”是两件事，各有一个条件。守恒的条件是规律本身不带时钟；而哈密顿量等于通常意义的总能量，还要求系统不受随时间变化的约束。比如一台被人从外面按着时间表推动的秋千，它的规律里写着时刻，能量可以不守恒——但这恰恰是时间平移对称被外力破坏的情形，不是反例，反而是定理的又一次正确预报。

\begin{mathbridge}
拉格朗日语言里，平移对称的推导只有两行。设系统由坐标 \(q\) 描述，拉格朗日量 \(L=T-V\) 不显含 \(q\)——这正是“规律与位置无关”的写法。欧拉—拉格朗日方程给出
\[
\frac{d}{dt}\frac{\partial L}{\partial \dot q}=\frac{\partial L}{\partial q}=0
\quad\Longrightarrow\quad
p=\frac{\partial L}{\partial \dot q}=\text{常量},
\]
这个 \(p\) 就是与 \(q\) 对应的动量。\(q\) 因此得到一个形象的名字：循环坐标——规律里找不到它，它对应的动量就守恒。时间平移同样一行：哈密顿量随时间的总变化等于它对时间的偏导数，
\[
\frac{dH}{dt}=\frac{\partial H}{\partial t},
\]
所以 \(H\) 不显含 \(t\) 时，\(H\) 本身就是守恒量。诺特论证的一般形状也在这里了：若某个连续改动 \(q\to q+\varepsilon\) 让 \(\delta L=0\)，那么沿真实路径，\(\partial L/\partial \dot q\) 沿这个方向的分量保持不变。对称改动的“方向”，决定守恒量的“种类”。
\end{mathbridge}

\begin{challenge}
诺特定理的一般形式：作用量 \(S=\int L\,dt\) 在某个连续变换下不变（\(\delta S=0\)），就存在一个沿真实运动守恒的电荷式数量 \(Q\)。把同样的思想用到场上，结论升级成一条局域定律：对称给出一个守恒流，满足连续性方程 \(\partial\rho/\partial t+\nabla\!\cdot\!\mathbf j=0\)——任何体积内“密度”的变化只能归因于它流过了边界，不能凭空生灭。电荷守恒就是这类局域守恒的最重要例子，它背后的对称不再是对时间和空间的朴素改动，而是“描述中的人为选择”——这正是第~59~章要讲的规范对称。还有一个让初学者不安的角落值得提前点名：在膨胀的宇宙里，光的波长一路上被拉长，光子能量一路变小，“宇宙的总能量守恒吗”不能用中学账本回答，因为膨胀的宇宙没有严格的时间平移对称——不是账本被撕了，而是“时间平移”这一页被宇宙的膨胀划掉了。
\end{challenge}

\begin{figure}[htbp]
\centering
\begin{tikzpicture}[>=Stealth]
  % 伸臂
  \draw[thick] (0,0) circle (1.5);
  \fill (0,0) circle (0.10);
  \fill (1.5,0) circle (0.14);
  \fill (-1.5,0) circle (0.14);
  \draw[->,thick] (45:1.95) arc (45:120:1.95);
  \node at (0,-2.15) {伸臂：\(I\) 大，\(\omega\) 小};
  \node at (0,2.35) {\(L=I\omega\)};
  % 收臂
  \draw[thick] (6,0) circle (0.6);
  \fill (6,0) circle (0.10);
  \fill (6.6,0) circle (0.14);
  \fill (5.4,0) circle (0.14);
  \draw[->,thick] ($(6,0)+(45:1.35)$) arc (45:160:1.35);
  \node at (6,-2.15) {收臂：\(I\) 小，\(\omega\) 大};
  \node at (6,2.35) {\(L=I\omega\)（同一个数）};
  \draw[->,very thick,dashed] (2.3,0) -- (3.7,0);
\end{tikzpicture}
\caption{收臂前后：质量分布收拢，转动惯量 \(I\) 缩小，角速度 \(\omega\) 按反比增大，乘积 \(L=I\omega\) 不变。动能不守恒——多出来的转动动能来自收臂时肌肉做的功。}
\end{figure}

\step{模型的边界}

第一条边界：对称的是规律，不是状态。地球表面附近，“上”和“下”是有特权方向的——引力替你选了。所以旋转对称只剩下绕竖直轴的那一份，花滑运动员绕竖直轴的角动量守恒，而绕水平轴翻跟头时重力矩随时可以改变账面上的数。状态可以不对称，规律仍然对称；方程对称而具体处境不对称，会发生什么有趣的事，是第~60~章的专题。

第二条边界是一个会让直觉翻车的反例：落猫。猫从高处四脚朝天落下，能在空中翻过半圈、四脚稳稳着地——全过程总角动量始终是零。它违反角动量守恒了吗？没有。猫不是刚体：它先把腰折成两截，让前半身和后半身朝相反的方向拧，再配合伸腿收腿改变两截的转动惯量，两截身体的角动量此消彼长、总和恒为零，整体姿态却换了个朝向。角动量守恒限制的是总量，从不限制内部各部分的重新分配。工程师照搬了这只猫：卫星和空间站肚子里的反作用轮——电机让飞轮正转，星体就反转；三个方向的飞轮配合，不喷一丝气体就能掉头。哈勃望远镜几十年如一日地稳稳指向深空，靠的就是这笔不花钱的账。

第三条边界是失效模式之一：账本的边界画小了。地面上推箱子，箱子的动量没了——它流进了地球；刹车时动能没了——它变成了刹车片和轮胎的内能（第~9~章和第~28~章细说过这笔账）。账永远平，只是科目可能散在别处。

第四条边界是失效模式之二：对称本身被破坏了。在加速前进的火车车厢里做碰撞实验，动量看起来不守恒——车厢不是一个平移对称的环境，外部世界正通过车轮对它使力。更深刻的例子在宇宙尺度：宇宙在膨胀，遥远星系的光在路上波长被不断拉长，光子能量不断变小。这不是能量守恒被推翻了，而是时间平移对称在膨胀宇宙里严格说来不成立——诺特定理在这里反而最好用：它明确告诉你哪一页账本失效了、为什么失效。

最后，用一组真实数据看看这条账本在自然界里记得有多认真。先给地球的自转记一笔总账：把地球近似成转动惯量约 \(8\times10^{37}\)~kg\(\cdot\)m\(^2\) 的球，自转角速度约 \(7.3\times10^{-5}\)~rad/s，乘起来，地球自转角动量约 \(5.9\times10^{33}\)~kg\(\cdot\)m\(^2\)/s——这是个庞大到没有日常类比的数量。可潮汐摩擦像刹车一样，一刻不停地拖着它：一天的长度大约每一百年变长 \(1.7\)~毫秒。地球自转角动量在减小——它去了哪？答案是月球轨道。阿波罗宇航员留在月面的反射镜让激光测距精确到厘米：月球正以每年约 \(3.8\)~厘米的速度远离地球，轨道角动量相应增加，增加的正是地球自转损失的那份。几亿年前，一天只有约二十二小时，这是珊瑚化石的生长纹记录下来的。角动量守恒不是教科书里的游戏，它在地质时间尺度上一笔一笔地记账，两端分毫不差。

\step{回头解释开场现象}

现在把开场的三个悬念一次性结清。

运动员的收银机。冰面附近，物理规律绕竖直轴旋转对称，所以绕竖直轴的角动量 \(L=I\omega\) 必须保持不变。收臂把质量拉向转轴，\(I\) 缩小，\(\omega\) 只能精确地按反比增大——“收银机”不只是比喻，它背后是几何：规律对朝向无所谓，转动账就必须平。她变快也不需要推力：增加的是转动动能而不是角动量，动能由肌肉从化学能的账上划拨而来。你在转椅上歪身子时那股别扭的晃荡，则是同一原理的反面教材：引力选定了“下”，旋转对称只剩竖直轴一份，绕其他方向的转动账不再受到保护，身体立刻感觉得到。

跳船的对称账目。水平方向上，人和船组成的系统近似孤立；空间平移对称保证相互作用成对抵消，总动量从零保持为零。你向前的动量有多大，船向后的就有多大——质量乘速度必然相等且反向，这不是巧合，是“规律处处相同”的直接推论。

实验隔夜不变。能量守恒等价于“物理规律不随时间变化”。你放心沿用昨天的仪器常数，是因为宇宙没有时间上的特价时段。反过来，如果有一天你发现秋千今天的周期和昨天不一样，你要排查的“故障”就是时间平移对称被破坏了——比如温度变了，比如发条松了。

至此，“先猜一猜”的三题也有了答案：守恒律不是独立的禁令，而是对称的影子（第一题选乙）；三次搬家里只要规律不变，结果就不变，而结果若变，对应的守恒律就地失效（第二题）；“热”不是遮羞布，是可以独立测量的内能——它恰恰是把时间平移对称的账算到底之后，必须存在的那一格科目（第三题）。禁令消失了，结构留下了。第~57~章问“什么叫对称”，本章回答了“对称会干什么”；第~59~章将继续追问：电势的零点、波函数的相位这类“描述中的人为选择”也是一种对称，而它对应的守恒量与相互作用本身有关——对称不只是管账，它还会立法。

\step{脑洞实验与挑战题}

\begin{thought}[深空中的封闭实验室]
设想把一间设备齐全的实验室装进没有窗的舱室，送到远离一切星体的深空，匀速漂移。你在舱内做任何力学实验：抛球、碰撞、摆钟、转椅。你能判断实验室“在哪里”“朝哪个方向”“现在是什么时候”吗？诺特定理的回答是：不能——这三个“不能”，正是动量、角动量、能量守恒的另一张面孔。但有一件事你能察觉：让舱室转起来，你不看窗外也知道自己在转——舱里的水会凹成抛物面，摆会偏出方向。转动似乎比位置、方向、时刻“更绝对”。为什么加速和转动能被察觉，而位置不能？这个问题困扰过马赫和爱因斯坦，一直通向广义相对论的深处；本书只把门推开一条缝。
\end{thought}

\begin{puzzles}
\item 一名宇航员出舱作业时安全绳断了，相对飞船静止地飘在三米之外。他身上没有推进器，手里只有一把两公斤的扳手。他怎样才能回到飞船？如果他不缺时间、只想转过身去看看地球，不扔任何东西能做到吗？\\
\emph{思路提示}：回飞船是“平移”问题——总动量必须保持为零，所以他必须向后扔出点什么；转身是“转动”问题——想想那只猫，总量可以为零而姿态可以改变，前提是身体（或手里的扳手）能分截拧动。

\item 假想一个“规律随地点缓慢变化”的宇宙角落：两球互推时，作用与反作用不再等大。那里会出现什么宏观怪象？再假想“规律随昼夜交替变化”：利用本章的奸商论证，设计一台在能量上净赚的装置。\\
\emph{思路提示}：对称破缺意味着对应账本的某一页作废——先找出作废的是哪一页，免费午餐就藏在那里。孤立物体会不会“自己动起来”，是第一个角落该出现的怪象。

\item 圆环形空间站靠自转制造“人工重力”。一名宇航员沿转动方向奔跑，另一名逆着转动方向奔跑，谁脚下的“体重秤”读数更大？两人跑起来时，空间站本身的自转会怎样？\\
\emph{思路提示}：把“人加空间站”当成一个孤立系统，总角动量不变。顺着跑的人带走的角动量变多了，舱体剩下的就变少，转速微微下降；脚下的“重力”来自旋转，转得快慢变了，读数自然一边偏大一边偏小。

\item 膨胀宇宙中，一束光从一百亿光年外出发，抵达地球时波长被拉长了一倍，每个光子的能量只剩一半。这些“丢失”的能量去了哪里？\\
\emph{思路提示}：先别急着找“能量藏到哪”，先检查前提——膨胀的宇宙里，时间平移对称还严格成立吗？账本失效的那一页，叫什么名字？
\end{puzzles}
